\clearpage
\thispagestyle{empty}


\begin{center}
	
	% (optional) delete marker line like your screenshot
	% {\small (English summary) $\leftarrow$ delete\par}
	% \vspace{0.5cm}
	
	% Title
	{\fontsize{20pt}{24pt}\selectfont
		\ThesisTitleKr \par}
	
	\vspace{1.6cm}
	
	% Author
	{\fontsize{14pt}{16pt}\selectfont
		\ThesisAuthorKr \par}
	
	\vspace{0.4cm}
	
	% Affiliation block (centered, stacked)
	{\fontsize{14pt}{24pt}
		\selectfont
		\ThesisSchoolKr \par
		\ThesisUniversityKr \par
	}
	
	\vspace{1.0cm}
	
	% Abstract heading
	{\fontsize{14pt}{16pt}\selectfont 요  \hspace{1cm} 약 \par}
	
\end{center}


% Body text (normal paragraph formatting)
\noindent
본 연구에서는 대표적인 천연 고분자의 하나인 전분을 여러 가소제를 사용
하여 가소화 시킨 thermoplastis starch(TPS)를 제조하여 이를 지방족 폴리에
스터의 하나인 polybutylene succinte(PBS)와 혼합비를 달리하며 블렌드 하
였다. TPS의 조성 및 함량이 PBS의 열적 특성 및 물리적 특성에 미치는 영향
에 대해 확인하였다. 
본 연구에서는 대표적인 천연 고분자의 하나인 전분을 여러 가소제를 사용
하여 가소화 시킨 thermoplastis starch(TPS)를 제조하여 이를 지방족 폴리에
스터의 하나인 polybutylene succinte(PBS)와 혼합비를 달리하며 블렌드 하
였다. TPS의 조성 및 함량이 PBS의 열적 특성 및 물리적 특성에 미치는 영향
에 대해 확인하였다. 
fillers.

% Keywords at bottom as an unnumbered footnote
\blfootnote{Keywords: thermoplastic starch (TPS); poly(butylene succinate) (PBS); biodegradable blends}
